\documentclass[10pt,a4paper]{amsart}
\usepackage{amsmath,amssymb,amsthm,graphicx}
\usepackage{amsmath}
\usepackage{amssymb}
\usepackage{hyperref}
\renewcommand{\theenumi}{\Alph{enumi}}
\renewcommand{\thesubsection}{\thesection\Alph{subsection}}

%%%%%%%%%%%%%%%%%%%%%%%%%%%%%%%%%%%%%%%%%%
% ENVIRONMENTS
%%%%%%%%%%%%%%%%%%%%%%%%%%%%%%%%%%%%%%%%%%%
\newtheorem{theorem}{Theorem}[section]
\newtheorem{prop}[theorem]{Proposition}
\newtheorem{lemma}[theorem]{Lemma}
\newtheorem{remark}[theorem]{Remark}
\newtheorem{corollary}[theorem]{Corollary}
\newtheorem{definition}[theorem]{Definition}
\newtheorem{example}[theorem]{Example}
%
\def\beq{\begin{equation}}
\def\eeq{\end{equation}}
\def\be{\begin{equation*}}
\def\ee{\end{equation*}}
%%%%%%%%%%%%%%%%%%%%%%%%%
% Beginning of the text %
%%%%%%%%%%%%%%%%%%%%%%%%%
%\begin{document}



%\documentclass[10pt]{article}
\usepackage{amssymb}
\usepackage{amsfonts}
%\topmargin 0.3in
%\textwidth 6in
%\textheight 9in%\baselineskip
%\hoffset=-50pt
%\voffset=-2pt
%\newenvironment{proof}[1][Proof]
%{\ \rule{0.5em}{0.5em}}


\title{Dedekind-finite Cardinals having countable partitions}
\author{Supakun Panasawatwong and J K Truss} 
\date{This paper is an amplified version of part of the first author's PhD thesis at the University of Leeds}%\footnotemark}
\title{Distinct infinite codes force the tree sets to be definably inequivalent inside Delta\_5}
\author{Supakun Panasawatwong, J. K. Truss}
\date{arxiv:2204.05700v1; CC BY 4.0}
\begin{document}
\maketitle
\noindent\textit{Editorial note.} All delivered mathematical text is the author's own, copied from the pinned publisher source sources/dedekind-finite-trees-2204.05700v1/source0.tex; no statement or proof step was rewritten, reordered or removed.\par\medskip
\noindent\textit{Editorial note.} Preamble provenance: the retained preamble is the source's own layout adapter, if any, plus the author's macro block at lines 5-44 of sources/dedekind-finite-trees-2204.05700v1/source0.tex (everything between the publisher class metadata and \textbackslash{}begin\{document\}, comments excluded). Nothing was added, no macro was invented, and the publisher document class is deliberately not reproduced; the excerpt serves only as preamble, never as a span coordinate.\par\medskip
\noindent\textit{Editorial note.} Metadata provenance: source\_id and source\_version were taken from focused-seeds/logic/STK-logic-code-inequivalence/license-reuse.json; the title from the card spec, taken from the publisher abstract page meta; the author list from the card spec, taken from the publisher abstract page meta (citation\_author); the source URL from the first non-empty of source\_url, source\_urls[0] and the evidence's fixed\_version\_url/license\_url. The licence evidence file used by this card is focused-seeds/logic/STK-logic-code-inequivalence/license-reuse.json, recorded in license.evidence\_path so that the binding target is explicit.\par\medskip
\noindent\textit{Editorial note.} Licence binding: the licence evidence staged with this source (LICENSE\_EVIDENCE.json) carries no status field and none of the accepted notice forms, so this card points at a per-card reuse record written inside its own task directory (license-reuse.json) with the fields name / version\_url / prior\_license\_path / basis; it asserts a material basis only and makes no new licence claim.\par\medskip
\noindent\textit{Editorial note.} External references are navigation pointers only; every deliberately kept step of the author's argument stays inside the counted spans. The displayed numbers of those pointers are the author's own numbering under the author's \textbackslash{}newtheorem declarations (the source has no subsections, so numbering is section.index), computed mechanically by the lane's scanner.\par\medskip
\noindent\textit{Editorial note.} The counted proof contains no reference to any labelled item and no bibliography citation, so the card needs neither a navigation stub nor a bibliography excerpt.\par\medskip
\noindent\textit{Editorial note.} Metadata provenance: identifiers come from the staged licence evidence of this source directory; the title and author list come from the citation\_title and citation\_author meta entries of the abstract page arxiv-abs.html kept beside it, as for the other cards of this source.\par\medskip
\noindent\textit{Allowed standard prerequisites (editorial).}\begin{itemize}
\item The definable cardinality classes Delta\_4 and Delta\_5 and the relation of definable equivalence between two sets in their respective permutation models. Set up earlier in the paper and referred to in the counted statement. Fixes the comparison that is refuted and the conclusion that both sets lie in Delta\_5 but not in Delta\_4.
\item The code of a weakly 2-transitive tree, recording the ramification orders along a maximal chain. Defined in the paragraph preceding the statement (source line 375). Supplies the invariant that the proof recovers from an assumed definable equivalence.
\item Weak 2-transitivity and density of the maximal chains of T\textasciicircum{}+, giving automorphisms that move points inside the pieces U\_a, L\_ab and C\_x while fixing a finite support. Used inside the counted proof at lines 382, 400-407 and 435-446. Lets the proof force incomparabilities and agreements of the two orders piecewise from the invariance of the support.
\end{itemize}\medskip
\begin{theorem}  \label{3.2} For any countable weakly $2$-transitive trees $T_1$, $T_2$ in which all maximal chains
of $T_1^+$, $T_2^+$ are dense, having distinct infinite codes, $|U_{T_1}|$ and $|U_{T_2}|$ are $\equiv$-inequivalent 
members of $\Delta_5 \setminus \Delta_4$.   \end{theorem}

\begin{proof}  The point of the restriction we have made on our weakly 2-transitive trees is that for any
ramification point $r$ and $x, y > r$ in $T$, there is an automorphism of $T$ fixing $r$ and taking $x$ to $y$.

Suppose for a contradiction that $|U_{T_1}| \equiv |U_{T_2}|$. In ${\mathfrak N}_{T_1}$ there must therefore be a 
structure $\prec$ on $U_{T_1}$ such that $(U_{T_1}, \prec)$ in ${\mathfrak N}_{T_1}$ is elementarily equivalent to 
$(U_{T_2}, <)$ in ${\mathfrak N}_{T_2}$. Let $A$ be a finite support for $(U_{T_1}, \prec)$ in ${\mathfrak N}_{T_1}$. 
Thus any automorphism of $(U_{T_1}, <)$ that fixes $A$ pointwise also fixes $\prec$. We may assume that $A$ is a 
subtree of $U_{T_1}$.

In order to handle this situation, we use the following notation, where $A$ is a finite subset of $T$ and $[A]$ 
is its definable closure, in this case just the lower subsemilattice generated by $A$. If $a$ is maximal in $A$ 
(or if $A = \emptyset$, $a$ can be $-\infty$), we let $U_a = \{x \in T: a < x\}$, and if $a < b$ are consecutive 
members of $[A]$, then $L_{ab} = \{x \in T: a < x < b\}$ (the linear piece between $a$ and $b$; again here $a$ 
can be $-\infty$ in the case when $b$ is minimal in $A$), and $S_{ab} = \{x \in T: a < x, b \not \le x\}$ (the 
corresponding side piece). Now $L_{ab}$ is linear, and we may also consider $L_{ab}^+$, which is the same thing, 
but calculated in $T^+$, that is, including ramification points too, which are coloured according to which orbit 
they lie in. If $x \in L_{ab}^+$, then $C_x = \{y \in T: b \wedge y = x\}$, which is the set of points which 
branch off from $L_{ab}^+$ at $x$. This is a union of cones at $x$.

A key remark is that the pointwise stabilizer $G_A$ acts transitively on each $U_a$, $L_{ab}$, and $C_x$ (here heavily 
using the assumption that the maximal chains in $T^+$ are dense). But more is true. 
In fact for any $b, c \in U_a$, there is an automorphism of $T$ taking $b$ to $c$ which fixes $A$ pointwise, {\em and} has
support contained in $U_a$, with a similar statement holding for $C_x$. The analogous statement also holds for
weak 2-transitivity. For instance, if $a < b < c$ and $a < d < e$, then by weak 2-transitivity there is an 
automorphism $f$ fixing $a$ and taking $b$ to $d$, and applying weak 2-transitivity again, there is an automorphism
$g$ fixing $d$ and taking $f(c)$ to $e$, and $f$, $g$ may be chosen fixing all points below $a$ and $d$ respectively.
Then $gf$ fixes $a$ and takes $b$ to $d$ and $c$ to $e$.

Although not actually needed, we can describe the orbits of $G_A$. They are of the form $U_a$, and $L_{ab}$, as well 
as the union of the cones in $S_{ab}$ which meet $L_{ab}$ at points lying in some orbit of $G_A$ on $L_{ab}^+$.

Now being a tree with dense chains is first order expressible, and so this must be true of $(U_{T_1}, \prec)$. We shall show
that on each set of the form $U_a$ or $C_x$, $\prec$ is equal to $<$ or the empty set (in which case it would be an antichain, 
i.e. pairwise $\prec$-incomparable). We just treat $U_a$, as $C_x$ is similar.

First consider a maximal point $a$ of $A$ (and let $a$ have a default value of $-\infty$ if $A = \emptyset$). First we note
that if $b, c \in U_a$ are $<$-incomparable, then they are also $\prec$-incomparable. This is because there is an 
element of $G_A$ swapping $b$ and $c$. Since $G_A$ preserves $\prec$, $b \prec c \Leftrightarrow c \prec b$, and since
$\prec$ is antisymmetric, $\neg b \prec c \wedge \neg c \prec b$. 

Next we show that there cannot be points $b, c \in U_a$ such that $b < c$ and $c \prec b$. For, if so, by weak
2-transitivity of $G_A$ on $U_a$, for $x, y \in U_a$, $x < y \Rightarrow y \prec x$. Choose $z > a$ and incomparable 
$x, y > z$. Thus $x, y \prec z$. By the previous paragraph, $x$ and $y$ are also $\prec$-incomparable, so this 
contradicts $(U_{T_1}, \prec)$ a tree. 

Next suppose that there are $b, c \in U_a$ such that $b < c$, and $b$ and $c$ are $\prec$-incomparable. By 
weak 2-transitivity of $U_{T_1}$, $a < x < y$ implies that $x$, $y$ are $\prec$-incomparable. But also as shown above, if
$b$ and $c$ are $<$-incomparable, they are also $\prec$-incomparable, and we deduce that $\prec$ is the empty relation
on $U_a$.

Otherwise, it follows that for $x, y \in U_a$, $x < y \Rightarrow x \prec y$. Since also $x, y$ $<$-incomparable
$\Rightarrow x, y \prec$-incomparable and $y < x \Rightarrow y \prec x$, we deduce that $x < y \Leftrightarrow x \prec y$. 
In other words, $<$ and $\prec$ agree on $U_a$.

Now let us consider $S_{ab}$ where $a < b$ are consecutive members of $A$. First note that by a proof similar to the above,
if $x, y$ are $<$-incomparable members of $C_c$ for $a < c < b$, then they are also $\prec$-incomparable, since there is a 
member of $G_A$ which swaps $x$ and $y$ (and fixes $c$). We can deduce that if $a < c < d < b$, and $x \in C_c$, $y \in C_d$,
then $x$ and $y$ are $\prec$-incomparable. For suppose that $x \prec y$ (a similar argument applying if $y \prec x$). Choose
$x' \in C_c$ which is $<$-incomparable with $x$. Then by what we have just shown, $x$, $x'$ are $\prec$-incomparable. 
Furthermore, there is a member of $G_A$ which swaps $x$ and $x'$, and having support contained in $C_c$, and so this also
fixes $y$. Since $G_A$ preserves $\prec$, $x' \prec y$, but this contradicts $\prec$ a tree.

It follows that any $\prec$-comparabilities in $S_{ab}$ must hold between members of the same $C_c$. Suppose then that
$x < y$ and $x \prec y$ hold for some $x, y \in C_c$. As $G_A$ acts weakly 2-transitively on $C_c$, $x < y \Rightarrow x \prec y$. 
Since we have also shown that $x$, $y$ $<$-incomparable implies that they are $\prec$-incomparable, we deduce that $<$ and $\prec$
agree on $C_c$. 

The main remaining point is to show that there is some $a$ or $c$ such that $\prec$ is non-empty on $U_a$ or $C_c$. Suppose not. Then
each $U_a$ is an antichain, and since the members of distinct $C_c$s for $a < c < b$ are incomparable, each $S_{ab}$ is also an
antichain. Consider the possible relationship between $a \in A$ and members of $U_b$ under $\prec$. The group $G_A$ acts transitively 
on $U_b$, and so if any member of $U_b$ is $\prec$-less than $a$, they all are, which violates $\prec$ a tree. Similarly, for members
of $C_c$. The same argument applies to the possible relationship between members of distinct $U_a$s, or $C_c$s, or between members 
of $U_a$ and $C_c$. We deduce that the union of all the $U_a$s and $S_{ab}$s is an antichain, and all its members are either 
incomparable with every member of $A$, or above some member of $a$. However, since $(U_{T_1}, \prec)$ is a tree satisfying the
same first order sentences as $(U_{T_1}, <)$, all its chains are infinite, so this is impossible.

We deduce that there is $a$ or $c$ such that $<$ and $\prec$ agree on $U_a$ or $C_c$. We just treat the first case, and show that 
for any $b \in U_a$, for $x \in T_1$, $b < x \Leftrightarrow b \prec x$. From left to right is already known. Suppose for a 
contradiction that there are some $b \in U_a$ and $x$ such that $b \prec x$ but not $b < x$. Since $<$ and $\prec$ agree on $U_a$, 
$x \not \in U_a$. Choose $c \in U_a$ incomparable with $b$. Then there is a member of $G_A$ taking $b$ to $c$ and with support 
contained in $U_a$. This preserves $\prec$, and fixes $x$, and hence $c \prec x$. Since $\prec$ is a tree, $b$ and $c$ are 
$\prec$-comparable, and as they both lie in $U_a$, they are also $<$-comparable, which gives a contradiction.

Finally we choose any $b \in U_a$, and observe that the ramification orders arising above $b$ in $<$ and $\prec$ are equal to those
arising in $T_1$ and $T_2$ respectively. Since these are determined by first order formulae, it follows that $|U_{T_1}|$ and $|U_{T_2}|$ 
have equal codes after all.  \end{proof}

\end{document}
